% from ega0-19b-fr.tex
\oldpage[0\textsubscript{IV-1}]{111}

\begin{definition}[19.8.4]
\label{0.19.8.4-fr}
Soient $A$ un anneau local, $P\to A$ l'unique homomorphisme d'un anneau local
premier $P$ dans $A$, $p$ la caractéristique des corps résiduels de $P$ et
$A$. On dit que $A$ est un \emph{anneau de Cohen} si c'est une $P$-algèbre de
Cohen, c'est-à-dire (19.8.1) si~:
\begin{enumerate}
  \item[$1^{\mathrm o}$] $A$ est noethérien et complet.
  \item[$2^{\mathrm o}$] $A$ est un $P$-module plat (ce qui revient aussi à
  dire que $A$ est un $\widehat P$-module plat (Bourbaki, \emph{Alg. comm.},
  chap.~III, §~5, n\textsuperscript{o}~4, prop.~4)).
  \item[$3^{\mathrm o}$] $A/pA$ est un corps (nécessairement séparable sur le
  corps résiduel de $P$, ce corps étant premier).
\end{enumerate}
\end{definition}

Si $p=0$, ces conditions équivalent à dire que $A$ est un \emph{corps de
caractéristique $0$}. Si $p>0$, on a nécessairement $pA\neq0$~; la condition
$3^{\mathrm o}$ signifie que $pA$ est l'idéal maximal $\mathfrak m$ de $A$~;
la condition $2^{\mathrm o}$ signifie que $p$ est $A$-régulier, puisque $P$
est un anneau de valuation discrète ($\mathbf 0_{\mathrm I}$, 6.3.4). Donc $A$
est un anneau régulier (17.1.1, $d)$) de dimension $1$, et par suite un anneau
de valuation discrète, complet en vertu de $1^{\mathrm o}$~; en résumé~:

\begin{proposition}[19.8.5]
\label{0.19.8.5-fr}
Les anneaux de Cohen sont les corps de caractéristique $0$ et les anneaux de
valuation discrète complets, dont le corps résiduel a une caractéristique
$p>0$, et dont l'idéal maximal est engendré par $p\cdot1$ ($1$ étant l'unité de
l'anneau).
\end{proposition}

On notera que dans le second cas, $p\cdot1\neq0$ puisque $p$ est $A$-régulier,
donc on peut identifier $p\cdot1$ à l'entier $p$~; l'homomorphisme canonique
$\mathbf Z_p\to A$ est injectif, et on identifie $p\cdot1$ à l'élément $p$ de
$\mathbf Z_p$~; on dit dans ce cas que $A$ est un \emph{$p$-anneau de Cohen}.

\begin{theorem}[19.8.6]
\label{0.19.8.6-fr}
\textup{(Cohen).} ---
\begin{enumerate}
  \item[(i)] Soient $W$ un anneau de Cohen, $C$ un anneau local noethérien
  complet, $\mathfrak J$ un idéal de $C$ distinct de $C$. Alors tout
  homomorphisme local $u:W\to C/\mathfrak J$ se factorise en
  $W\xrightarrow{v}C\to C/\mathfrak J$, où $v$ est un homomorphisme local.
  \item[(ii)] Soit $K$ un corps. Il existe un anneau de Cohen $W$ dont le corps
  résiduel est isomorphe à $K$. Si $W'$ est un second anneau de Cohen, $K'$ son
  corps résiduel, tout isomorphisme $u:K\xrightarrow{\sim}K'$ provient par
  passage aux quotients d'un isomorphisme $v:W\xrightarrow{\sim}W'$.
\end{enumerate}
\end{theorem}

\begin{proof}
Cela n'est autre que (19.8.2) appliqué au cas où $A$ est un anneau local
premier.
\end{proof}

\begin{remarks}[19.8.7]
\label{0.19.8.7-fr}
\begin{enumerate}
  \item[(i)] Lorsque $K$ est de caractéristique $0$, la partie (ii) de (19.8.6)
  devient triviale.
  \item[(ii)] L'homomorphisme $v$ de (19.8.6, (i)) \emph{n'est pas
  nécessairement déterminé de façon unique par $u$}, comme le montre déjà le
  cas où $W$ est un corps de caractéristique $0$, $\mathfrak J^2=0$ et $u$ est
  un isomorphisme (cf.~(21.5.5)). De même, dans (19.8.6, (ii)) l'isomorphisme
  $v$ \emph{n'est pas nécessairement déterminé de façon unique par $u$}
  (cf.~(21.5.5)).

  Toutefois, lorsque $K$ est \emph{parfait} et de caractéristique $p>0$, on
  verra (21.5.5) que dans (19.8.6, (ii)) l'isomorphisme $v$ est unique. On verra
  aussi plus tard que dans ce cas $W$ s'identifie à l'anneau $W_\infty(K)$ des
  \emph{vecteurs de Witt de longueur infinie} sur $K$.
  \item[(iii)] Dans (19.8.6, (i)), on peut affaiblir les hypothèses sur $C$ en
  utilisant (19.3.10) et (19.3.12).
\end{enumerate}
\end{remarks}

\begin{theorem}[19.8.8]
\label{0.19.8.8-fr}
\textup{(Cohen).} --- Soient $A$ un anneau local noethérien complet, $k$ son
corps résiduel.

